\documentclass[12pt]{article}
\usepackage{epsf, amssymb, amsmath, amsthm, hyperref, tikz, array, booktabs, calc, accents}
\usetikzlibrary{arrows, shapes, snakes, automata, positioning, backgrounds, petri, tikzmark, calc, decorations.pathreplacing}
\tikzset{every picture/.style={remember picture}}
\usepackage[english]{babel}
\newcommand\myubar[1]{%
\underaccent{\bar}{#1}}
\newtheorem{theorem}{Theorem}[section]
\newtheorem{lemma}[theorem]{Lemma}

\title{---}
\author{---}

\setlength{\oddsidemargin}{-0.5in}
\setlength{\topmargin}{-0.6in}
\setlength{\headheight}{0cm}
\setlength{\headsep}{0cm}
\setlength{\textheight}{10.3in}
\setlength{\textwidth}{7.5in}
\setlength{\topskip}{0cm}
\setlength{\parskip}{4mm}
\setlength{\parindent}{0cm}
\renewcommand{\baselinestretch}{1.15}
\renewcommand{\arraystretch}{1.15}
\renewcommand{\dbltopfraction}{0.9}
\renewcommand{\topfraction}{0.9}
\renewcommand{\bottomfraction}{0.9}
\renewcommand{\dblfloatpagefraction}{0.99}
\renewcommand{\floatpagefraction}{0.99}
\renewcommand{\textfraction}{0.01}


% numeric sets

\newcommand{\Reals}{\mathbb{R}}      % real numbers
\newcommand{\Naturals}{\mathbb{N}}   % natural numbers
\newcommand{\Integers}{\mathbb{Z}}   % integer numbers
\newcommand{\Rationals}{\mathbb{Q}}  % rational numbers
\newcommand{\Complexes}{\mathbb{C}}  % complex numbers

% math symbols

\newcommand{\IFF}{\mbox{$\Longleftrightarrow$}}    % biimplication
\newcommand{\THEN}{\mbox{$\Rightarrow$}}           % implication

\tikzset{
  node distance=1mm,
  terminal/.style={
    rectangle, % The shape:
    minimum size=1mm, % The size:
    thick, draw=black!75, % The border:
    fill=black!20, % The filling:
    font=\ttfamily
  },
  nonterminal/.style={
    align=center,
    circle, % The shape:
    minimum size=3mm, % The size:
    text width=3mm, % This forces nodes to have the same size
    thick, draw=black!75, % The border:
    fill=black!20, % The filling:
    font=\itshape,
    double, double distance=1mm % Makes the border thicker for final states
  },
}


\pagestyle{empty}

\begin{document}

%\maketitle

\noindent\begin{tabular}{p{4cm}p{4cm}p{10cm}}
\textbf{COMS 6730} & \textbf{Spring 2025} & \textbf{Name: \underline{Mobina Amrollahi\hspace{4cm}}}
\end{tabular}


\begin{center}
   \textbf{CheatSheet} 
\end{center}
\section*{Linear Models}
\begin{enumerate}
    \item \textbf{Linear regression} is a linear model, trying to predict a \textbf{continuous scalar}.
    \item \textbf{RSS} in linear regression minimizes the sum of the squared distances between the label of each data point and the predicted value.
    \item The weights can be found by \textbf{maximizing the likelihood of the data} under the assumption of Gaussian noise, which is equivalent to \textbf{minimizing the RSS}.
    \item Finding the weights by solving the normal equations might not be possible as $X^TX$ might not exist. 
    \item If $X$, the feature matrix, is \textbf{full rank}, then the optimal solution for linear regression exists and is unique.

    \item Three assumptions of linear regression: 1. Data is linear, 2. Data points are independent, 3. The residual follows a normal distribution with zero mean.
    \item The perceptron learning algorithm guarantees convergence for the dataset that is \textbf{linearly separable}, regardless of the initial data or weights.

    \item \textbf{Logistic regression} is a linear model because the decision boundary is a linear function of the input features.

    \item The \textbf{monotonicity} of the sigmoid function ensures that logistic regression produces a linear decision boundary. By changing sigmoid to sine for instance, logistic regression becomes a non-linear model.
\end{enumerate}

\section*{Gradient Descent Algorithm}
\begin{enumerate}
    \item Batch size affects both the \textbf{forward} and \textbf{backward} propagation.
    \item \begin{table}[h!]
        \centering
        \scriptsize
        \begin{tabular}{>{\raggedright}p{4cm}| >{\raggedright}p{4cm}| >{\raggedright\arraybackslash}p{4cm}}
        \toprule
        \textbf{Batch}
        & 
        \textbf{Mini - Batch} 
        & 
        \textbf{Stochastic} \\
        \midrule
        
        \begin{itemize}
            \item process the entire dataset at once
            \item smoother, more accurate gradient estimates
            \item computationally expensive in practice
        \end{itemize}
        
        & 
        
        \begin{itemize}
            \item helps escape local minima during optimization
            \item introduces noise into the optimization process 
            \item small to medium batch sizes can lead to better generalization
        \end{itemize}
        
        &
        
        \begin{itemize}
            \item update weights using one sample at a time
        \end{itemize}
        
        \\
        
        \bottomrule
        \end{tabular}
        \caption{Comparison between Different Algorithms}
    \end{table}
\end{enumerate}

    
\section*{Convolution}
    \begin{enumerate}
    \item Kernel is a small matrix. They're typically small and scan across the image.
    \item Feature maps: the result of applying filters to the input.
    \item The primary purpose of a convolution layer is to extract spatial features from the input.
    \item A convolution operation multiplies kernel and input element-wise, then sums them up.
    \item Hand-Crafted Kernels:
    \begin{table}[h!]
    \centering
    \small
    \begin{tabular}{>{\raggedright}p{4cm}| >{\raggedright}p{4cm}| >{\raggedright}p{4cm}| >{\raggedright\arraybackslash}p{4cm}}
    \toprule
    \textbf{Box Filter} & \textbf{Sharpening Filter} & \textbf{Sobel (Vertical Edge)} & \textbf{Sobel (Horizontal Edge)} \\
    \midrule
    % Filter matrices
    \[
    \begin{bmatrix}
    \frac{1}{9} & \frac{1}{9} & \frac{1}{9} \\
    \frac{1}{9} & \frac{1}{9} & \frac{1}{9} \\
    \frac{1}{9} & \frac{1}{9} & \frac{1}{9} \\
    \end{bmatrix}
    \] 
    & 
    \[
    \begin{bmatrix}
    0 & -1 & 0 \\
    -1 & 5 & -1 \\
    0 & -1 & 0 \\
    \end{bmatrix}
    \] 
    & 
    \[
    \begin{bmatrix}
    -1 & 0 & 1 \\
    -2 & 0 & 2 \\
    -1 & 0 & 1 \\
    \end{bmatrix}
    \]
    & 
    \[
    \begin{bmatrix}
    -1 & -2 & -1 \\
    0  & 0  & 0 \\
    1  & 2  & 1 \\
    \end{bmatrix}
    \] \\
    
    % Properties
    \midrule
    
    \begin{itemize}
        \item Averages neighboring pixels
        \item Smooths the image
        \item Reduces noise
        \item Removes sharp features
    \end{itemize}
    &
    \begin{itemize}
        \item Enhances edges
        \item Emphasizes differences with neighbors
        \item Increases sharpness
        \item High-pass filter
    \end{itemize}
    &
    \begin{itemize}
        \item Detects vertical edges
        \item Computes horizontal gradients
        \item Sensitive to vertical features
        \item Emphasizes vertical boundaries
    \end{itemize}
    &
    \begin{itemize}
        \item Detects horizontal edges
        \item Computes vertical gradients
        \item Sensitive to horizontal features
        \item Emphasizes horizontal boundaries
    \end{itemize}
    \\
    
    \bottomrule
    \end{tabular}
    \caption{Comparison of Different Filters and Their Properties}
    \end{table}
    \item Despite hand-crafted filters, the learned ones are randomly initialized, and then optimized through back-propagation.

    \item Convolution uses a set of kernels, each applied to an individual input channel.

    \item The main effect of \textbf{stride} is that it reduces the output feature map.

    \item The main effect of \textbf{padding} is to ensure all pixels are equally used.

    \item Number of trainable parameters in each convolution layer: $K \times K \times C_{in} \times C_{out}$.

    \item Number of multi-add operations or computational cost in each convolution layer: $K \times K \times C_{in} \times C_{out} \times W' \times D'$.

    \item \textbf{Dilated convolutions} enlarges the receptive field by introducing gaps (holes) in the kernel to cover a larger area.  For a dilation “rate” $d$, $d-1$ spaces are inserted between kernel elements such that $d = 1$ corresponds to a regular convolution.

    \item Output dimension: $W' = \frac{W-K+2P -(k-1)(d-1)}{S} + 1$.

    \item Receptive Field: the region of the input space that affects a particular unit in the network.
    \[\forall \quad l: R_l = R_{l-1} + (k_{l} - 1) \times J_{l-1} \times d_l\]
    The jump $J_l$ describes how far we move in the input when moving one unit in the feature map:
    \[
    J_{l} = J_{l-1} \times S_{l}, \quad \text{and} \quad J_0 = 1, \quad R_0 = 1
    \]
    where
    \begin{itemize}
        \item $R_l$: receptive field size at layer $l$
        \item $k_{l}$: kernel size at layer $l$
        \item $J_{l-1}$: jump (effective stride) from the previous layer
        \item $S_{l}$: stride at layer $l$
        \item $d_l$: dilation rate at layer $l$
    \end{itemize}
    In which, you can see for a network with one layer, only \textbf{the kernel size} can affect the receptive field. However, once the number of kernels increases, the \textbf{stride} and \textbf{pooling layers} also affect it.

    \item $1 \times 1$ Convolution: is used to reduce the number of channels (dimensionality) while introducing non-linearity.

    \item \textbf{Pooling layers} reduce the output feature maps, while avoid overfitting, and enlarging the receptive field.

    \item Solving the high dimension problem:
    \begin{equation*}
    \tikz[baseline]{\node(d1) {$W$}} \times \tikz[baseline]{\node(d2){$D$}} \times \tikz[baseline]{\node(d3) {$C$}}
    \end{equation*}
    \vspace{1cm}\par
    % insert text
    \noindent\hfil\begin{tikzpicture}
    \node (t1) {Increasing the Stride};
    \node[above right] (t2) at (t1.east) {Max/Avg Pooling Layers};
    \node[right,text width=2.5cm,align=center] (t3) at (t2.east)
     {$1 \times 1$ Convolution};
     \node[below left, align=center] (t4) at (t1.west)
     {Depthwise Separable Convolution};
    \end{tikzpicture}
    % insert lines between text
    \begin{tikzpicture}[overlay]
    \draw[blue,thick,->] (d1) to [in=90,out=245] (t1.north);
    \draw[green,thick,->] (d1) to [in=90,out=245] (t2.north);
    \draw[red,thick,->] (d1) to [in=90,out=245] (t4.north);
    \draw[blue,thick,->] (d2) to [in=90,out=265] (t1.north);
    \draw[green,thick,->] (d2) to [in=90,out=265] (t2.north);
    \draw[red,thick,->] (d2) to [in=90,out=245] (t4.north);
    \draw[black,thick,->] (d3) to [in=90,out=265] (t3.north);
    \end{tikzpicture}

    \item Activation functions are element-wise function that introduce non-linearity. Tanh and sigmoid have gradient vanishing problems, meanwhile ReLU has the problem of \textbf{dying neurons}. One attempt to fix the dying neurons problem is to use \textbf{leaky ReLU}.

    \item  \textbf{Nearest Neighbors}, \textbf{Bi-Linear Interpolation}, \textbf{Bed of Nails}, and \textbf{Max-Unpooling} are all \textbf{deterministic} up sampling techniques, meanwhile transposed convolution is learnable.

    \item \textbf{Bed of Nails} is followed by a convolution layer that can interpolate or blend the sparse data points to generate meaningful and smooth output.

    \item \textbf{Max-Unpooling} it's always following a corresponding max-pooling layer. 

    \item Transposed convolution is often \textbf{incorrectly} called \textbf{de-convolution}. However de-convolution refers to the inverse operation of standard convolution. The name \textbf{transposed} corresponds to multiplying by the transpose of the convolution kernel matrix. Transposed convolutions swap the forward and backward passes of a convolution.
    \textbf{Output size formula}:
    \[
    W_{out} = S \times (W - 1) + K - 2P
    \]

    \item In transposed convolution if $s > 1$ we will put $s -1$ zeros between input elements and in the borders.

    \item Unlike deep convolutional networks, \textbf{attention} mechanisms can capture long-range dependencies in a single layer. They're called global extractors, because each query attends all the keys and values.

    \item \textbf{Scaled Dot-Product Attention:}
    
    \[
    \text{Attention}(Q, K, V) = \text{Softmax}\left( \frac{QK^T}{\sqrt{d_k}} \right) V
    \]
    
    \begin{itemize}
        \item $Q \in \mathbb{R}^{n \times d}$  
        \item $K \in \mathbb{R}^{m \times d}$ 
        \item $V \in \mathbb{R}^{m \times p}$  
        \item $d_k$: Dimensionality of keys (used for scaling)
    \end{itemize}

    \item \textbf{Self-Attention with a Single Input Matrix $X$:}
    \[
    Q = X W_Q, \quad K = X W_K, \quad V = X W_V
    \]
    
    \[
    \text{Attention}(X) = \text{Softmax}\left( \frac{(X W_Q) (X W_K)^T}{\sqrt{d_k}} \right) (X W_V)
    \]
    
    \begin{itemize}
        \item When using self-attention, $X$ is projected into $Q, K, V$ using trainable weight matrices.
        \item This allows each token in $X$ to attend to others, learning dependencies across positions.
    \end{itemize}

    \item \textbf{Convolution} and \textbf{attention} layers can be interpreted as variations of fully-connected layers, since they're all matrix multiplications. However, pooling layers (e.g., max pooling, average pooling) perform non-parametric reduction operations and are not typically expressed as matrix multiplications.
    \end{enumerate}
    
\section*{Convolution Neural Network}
\begin{enumerate}
    \item Winners of ImageNet Comparison:
    \begin{table}[h!]
    \centering
    \scriptsize
    \begin{tabular}{>{\raggedright}p{4cm}| >{\raggedright}p{4cm}| >{\raggedright}p{4cm}| >{\raggedright\arraybackslash}p{4cm}}
    \toprule
    \textbf{AlexNet} & \textbf{VGG} & \textbf{ResNet} & \textbf{DenseNet} \\
    \midrule
    \begin{itemize}
        \item 8 layers
        \item ReLU Activation Function
        \item High number of kernels at each layer
        \item Different kernel sizes of $11 \times 11, 5 \times 5, 3 \times 3$.
        \item Introduced Dropout
        \item Trained with Data Augmentation
        \item Overlap pooling layer
    \end{itemize}
    &
    \begin{itemize}
        \item 19 Layers
        \item ReLU Activation Function
        \item Fixed pattern of conv + conv + pool at each layer
        \item Fixed kernel size of $3 \times 3$.
        \item Max pooling $2 \times 2$ with stride $=2$
        \item Same padding everywhere
        \item Has a degradation problem: vanishing gradients or explosion
        \item Prone to Overfit
    \end{itemize}
    &
    \begin{itemize}
        \item 152 Layers
        \item Skip connection: preventing the gradient vanishing 
        \item Batch Normalization: preventing the gradient explosion 
        \item Each residual network: $3 \times 3$ conv + Batch Normalization + ReLU + $3 \times 3$ conv + Batch Normalization
    \end{itemize}
    &
    \begin{itemize}
        \item Concatenates the outputs from different layers.
        \item Computes vertical gradients
        \item Uses more memory compare to ResNet
    \end{itemize}
    \\
    
    \bottomrule
    \end{tabular}
    \caption{Comparison of Different Deep Learning Architectures}
    \end{table}

    \item In \textbf{DenseNet}, the $l$-th layer receives as input the concatenation of all feature maps produced by previous layers:

    \[
    x_l = H_l([x_0, x_1, x_2, ..., x_{l-1}])
    \]
    
    Where:
    \begin{itemize}
        \item $[x_0, x_1, ..., x_{l-1}]$ represents the concatenation of feature maps along the channel dimension.
        \item $H_l(\cdot)$ is the composite function of operations: BatchNorm, ReLU, and Convolution.
        \item If the growth rate is $k$, the input depth to layer $l$ is:
        \[
        d_{in} = d_0 + k \times (l - 1)
        \]
    \end{itemize}

    \item In \textbf{Transfer Learning}:
    \begin{enumerate}
        \item Train a network like VGG or ResNet on a large dataset
        \item Freeze the earlier convolution layers weights
        \item Replace the fully connected layer with new layer specific to the task
    \end{enumerate}

    \item Weight Initialization in CNNs:
    \begin{table}[h!]
        \centering
        \scriptsize
        \begin{tabular}{>{\raggedright}p{4cm}| >{\raggedright}p{4cm}| >{\raggedright\arraybackslash}p{4cm}}
        \toprule
        \textbf{Xavier Initialization}
        & 
        \textbf{Kaiming Initialization} 
        & 
        \textbf{Random Initialization} \\
        \midrule
        
        \begin{itemize}
            \item $W \sim \mathcal{N} \left( 0, \frac{2}{n_{in} + n_{out}} \right)$
            \item Works well with sigmoid and tanh activations.
        \end{itemize}
        
        & 
        
        \begin{itemize}
            \item $W \sim \mathcal{N} \left( 0, \frac{2}{n_{in}} \right)$
            \item Best for ReLU and Leaky ReLU activations.
        \end{itemize}
        
        &
        
        \begin{itemize}
            \item $W \sim \mathcal{N}(0, 10^{-2})$
            \item Works on small networks
            \item still have problems on deep networks
            \item Weight gradients in deeper layers have variances of nearly zero
        \end{itemize}
        
        \\
        
        \bottomrule
        \end{tabular}
        \caption{Weight Initialization}
    \end{table}

    \item \textbf{Batch Normalization}: During training, the mean and variance are computed from each mini-batch. During testing, the moving average of the mean and variance, calculated during training, is used instead.

    \item \textbf{Batch size} has an impact in batch normalization.

    \item During the prepossessing we want the \textbf{input data} to be \textbf{normalized} and \textbf{zero-centered.}

    \item \textbf{Data augmentation} increases the diversity of the training dataset, forcing the model to generalize better. However, the labels of the augmented data remain the \textbf{same}.

    \item  During training, \textbf{Dropout} randomly deactivates neurons with probability $p$, which lowers the expected output. To maintain consistency between training and testing:

    \begin{itemize}
        \item \textbf{Inverted Dropout (common)}: During training, scale the output of active neurons by $\frac{1}{1 - p}$.
        \item \textbf{Standard Dropout}: During training, no scaling is applied. During testing, scale the activations by $(1 - p)$.
    \end{itemize}

    \item Since in the drop-out neurons will randomly be set to zero, the number of trainable parameters won't change.

    \item Drop-out is only applied to the \textbf{hidden} layer, and not the output layer.
\end{enumerate}

     

\end{document}
